\chapter{Basic Set Theory}
\label{part:vector-spaces-scalars-quotients}


\label{chap:foundations}

\section{Sets and Operations}
\label{sec:sets-basic-constructions}
A \textit{set} is just a collection of things, denoted (in common) by one letter, and its items are listed in the bracket $\{\}$ and separated by commas, for example:
\[
    A = \{1,2,3, \text{ orange}, \text{ apple}, \text{ anything you want}\}.
\]
The set of sets $(G = \{A,B\})$ is a set, and the set of nothing ($X = \{\varnothing\}$) is a set, but nothing ($\varnothing$) is itself a set, but $\varnothing$ is different to $\{\varnothing\}$ (why?). A \textit{subset $B$ of the set $A$} is also a set (of course), and every item in $B$ is contained in $A$, for example, this is an subset of $A$:
\[
    B = \{1,3, \text{ orange}\},
\]
then we write $A \subseteq B$. If we want to emphasize that ""$A$ contains $B$"", we want to write $A \supseteq B$. On the other hand, this is not a subset of $A$: 
\[
    C = \{2,3, \text{orange}, \text{ strawberry}\},
\]
and we write $C \subsetneq A$. Finally, we write $A = B$ if $A \subseteq B$ and $B \subseteq A$, most of proofs we would like to use this definition.

Now introduce some set operations: 
\begin{enumerate}
    \item $A \cap B$ denotes all items belonging to $A$ and $B$,
    \item $A \cup B$ denotes all items belonging to $A$ or $B$,
    \item $A \backslash B$ denotes all items belongis to $A$ but not belonging to $B$.
    \item $2^A$ denotes the set of all subsets of $A$.
\end{enumerate}  
Before the next section, try to prove some statements, it's easy but necessary.
\begin{enumerate}
    \item  If $A \cap B = \varnothing$, then $A \backslash B = A$.
    \item If $B \subseteq A$, then $A \backslash B = A \cap B$.
    \item 
\end{enumerate}

\section{Set Products}
\label{sec:ordered-pairs-products-relations}


Let $a$ and $b$ be items. The \textit{Kuratowski ordered pair} $(a,b)$ is defined by
\[
  (a,b):=\{\{a\},\{a,b\}\}.
\]
The first entry is called the \emph{first coordinate}, and the second
entry is called the \emph{second coordinate}.



\begin{proposition}
For any items \(a,b,x,y\), one has 
$
  (a,b)=(x,y)
$
if and only if
$a = x$ and $b = y$.
\end{proposition}

\begin{proof}
Assume first that $(a,b)=(x,y)$.
By the definition of an ordered pair, one has
\[
  \bigl\{\{a\},\{a,b\}\bigr\}
  =
  \bigl\{\{x\},\{x,y\}\bigr\}.
\]

Suppose first that $a=b$. Then
$
  (a,b)=\{\{a\}\}.
$ Hence $(x,y)$ must also contain only one
element, so $\{x\} = \{x,y\}$, implying that $x = y$ and $a = b = x = y$.

Now suppose that \(a\neq b\). Then
$
  \{a\}\neq\{a,b\},
$
so the set defining $(a,b)$ has two distinct elements. Since the singleton
$\{x\} \neq \{a,b\}$, then $\{x\}=\{a\}$,
and hence \(x=a\). It follows that
$
  \{x,y\}=\{a,b\}.
$
Since $x=a$ and $a \neq b$, then $y \neq x$ and $y = b$.

Conversely, if \(a=x\) and \(b=y\), then the two ordered pairs are equal
directly from the definition.
\end{proof}
Now we want to expand the ordered pair into the $n$-ordered tuple. Let $A$ and $B$ be sets. Their \textit{Cartesian product} is
\[
  A\times B
  :=
  \{(a,b):a\in A,\ b\in B\}.
\]
\begin{remark}
In general,
\[
  A\times B\neq B\times A.
\]
Even when the same objects occur in both products, their coordinates
appear in different positions.
\end{remark}

\begin{definition}[Cartesian powers]
For a set \(A\), define
\[
  A^2:=A\times A.
\]
More generally,
\[
  A^n
  :=
  \underbrace{A\times\cdots\times A}_{n\text{ factors}}.
\]
Equivalently, one may define the powers recursively by
\[
  A^{n+1}:=A^n\times A.
\]
\end{definition}

An element of \(A^n\) is written as
\[
  (a_1,\ldots,a_n),
\]
where every coordinate belongs to \(A\).

\begin{example}
If
\[
  A=\{1,\text{apple}\},
\]
then
\[
  A^2
  =
  \{
    (1,1),
    (1,\text{apple}),
    (\text{apple},1),
    (\text{apple},\text{apple})
  \}.
\]
\end{example}

\begin{proposition}
Let \(A\), \(B\), and \(C\) be sets. Then
\[
  A\times(B\cup C)
  =
  (A\times B)\cup(A\times C).
\]
\end{proposition}

\begin{proof}
Take an element \(z\in A\times(B\cup C)\). Then \(z=(a,t)\) for some
\(a\in A\) and some \(t\in B\cup C\). Hence \(t\) belongs to \(B\) or
to \(C\). In the first case, \(z\in A\times B\); in the second case,
\(z\in A\times C\). Therefore
\[
  z\in(A\times B)\cup(A\times C).
\]

For the reverse inclusion, take
\[
  z\in(A\times B)\cup(A\times C).
\]
If \(z\in A\times B\), then \(z=(a,b)\) with \(a\in A\) and
\(b\in B\subseteq B\cup C\). Hence \(z\in A\times(B\cup C)\).
The case \(z\in A\times C\) is identical.
\end{proof}

\begin{proposition}
Let \(A\), \(B\), and \(C\) be sets. Then
\[
  A\times(B\cap C)
  =
  (A\times B)\cap(A\times C).
\]
\end{proposition}

\begin{proof}
An ordered pair \((a,t)\) belongs to \(A\times(B\cap C)\) precisely
when \(a\in A\) and \(t\) belongs to both \(B\) and \(C\). This is
exactly the condition that \((a,t)\) belong to both \(A\times B\) and
\(A\times C\).
\end{proof}

\begin{proposition}
For every set \(A\),
\[
  A\times\varnothing=\varnothing
  \qquad\text{and}\qquad
  \varnothing\times A=\varnothing.
\]
\end{proposition}

\begin{proof}
An ordered pair in \(A\times\varnothing\) would need a second coordinate
belonging to the empty set, which is impossible. The proof for
\(\varnothing\times A\) is the same.
\end{proof}


\section{Relations}
\label{subsec:relations}

\begin{definition}[Relation from one set to another]
Let \(A\) and \(B\) be sets. A relation from \(A\) to \(B\) is a subset
\[
  R\subseteq A\times B.
\]
When \((a,b)\in R\), we write
\[
  a\,R\,b
\]
and say that \(a\) is related to \(b\) by \(R\).
\end{definition}

A relation is therefore a chosen collection of ordered pairs. It need
not assign exactly one element of \(B\) to each element of \(A\); that
additional requirement belongs to the definition of a function.

\begin{example}
Let \(A\) be a set of students and \(B\) a set of project topics. Define
\(R\subseteq A\times B\) by declaring that
\[
  s\,R\,t
\]
whenever the student \(s\) is assigned to the topic \(t\). Then \(R\)
is a relation from \(A\) to \(B\).
\end{example}

\begin{example}
Let
\[
  A=[-1,1],
  \qquad
  B=\R.
\]
The set
\[
  R
  =
  \{(x,y)\in A\times B:x^2+y^2=1\}
\]
is a relation from $A$ to $B$. Geometrically, it is the unit circle
viewed as a subset of the Cartesian plane.
\end{example}



